\section{A positive-label ReLU memorizer at nearly one unit per layer}
\label{sec:relu-near-baseline}
Use the unit-sphere support $\mathcal X_D$ of Section~\ref{sec:gaussian-feature-memorizer}, with $d=D+2$ and $m=4\cdot2^D$. Every augmented input has norm $\sqrt2$. Its uniform probability measure is understood in all inner products. Put $\Lambda_D^{\rm R}=2^{-16}e^{-12(D+4)}$.

\begin{lemma}[ReLU Gaussian features]
For $A\sim N(0,I_{d+1}/d)$ and $\phi_A(x)=\operatorname{ReLU}(A\cdot(1,x))$, the integral operator
$K(x,x')=\E_A\phi_A(x)\phi_A(x')$
has smallest eigenvalue at least $\Lambda_D^{\rm R}$.
Every sufficiently large integer $N$ admits $N$ such deterministic rows $a_i$ with
\[
 N^{-1}\sum_i\|a_i\|^2\le2,\qquad
 K_N=N^{-1}\sum_i\phi_{a_i}\otimes\phi_{a_i}\succeq\tfrac12\Lambda_D^{\rm R} I.
\]
It suffices that $N\ge C m^2(\Lambda_D^{\rm R})^{-2}$.
For any prescribed $r_0$ on the support there are coefficients satisfying
\[
 r_0=N^{-1}\sum_i c_i\phi_{a_i},\qquad
 N^{-1}\sum_i c_i^2\le2(\Lambda_D^{\rm R})^{-1}\|r_0\|_2^2.
\]
\end{lemma}
\begin{proof}
Expand $\operatorname{ReLU}(Z)$ in normalized probabilists' Hermite polynomials. For every even $j=2r\ge2$, its squared coefficient is
\[
 b_{2r}^2=
 \frac{\binom{2r-2}{r-1}}{4^{r-1}}\,
       \frac1{2\pi(2r)(2r-1)}
 \ge\frac1{2\pi(2r)(2r-1)^2}
 \ge\frac1{2\pi j^3}.
\]
The formula follows from Gaussian integration by parts:
$\E[\operatorname{ReLU}(Z)\operatorname{He}_{2r}(Z)]
=(-1)^{r-1}(2r-3)!!/\sqrt{2\pi}$.
Equivalently integrate the ordinary polynomial on the positive half-line twice and use the Gaussian derivative identity. The lower bound uses that the central binomial coefficient is the largest of $2r-1$ nonnegative coefficients summing to $4^{r-1}$.

The preactivation variance is $2/d$ at every input. By positive homogeneity,
\[
 K(x,x')=\frac2d\sum_{j\ge0}b_j^2
           \left(\frac{1+x\cdot x'}2\right)^j.
\]
For each nuisance subset $S$ and $a\in\{0,1,2,3\}$, choose $b\in\{0,1,2\}$ such that $j=|S|+a+b$ is even and at least two. One can arrange $j\le D+4$; use $b=2$ only when $|S|+a=0$. Retain the multinomial term
\[
 \frac2d\frac{b_j^2j!}{a!b!\,2^{a+b}(8D)^{|S|}}
       T^a\chi_S(z)(T')^a\chi_S(z').
\]
These selected terms are distinct and positive semidefinite. Their coefficients are at least $e^{-12(D+4)}$. To check this loose uniform bound, write $k=D+4\ge5$. Since $j\le k$, $a!b!\le12$, and $2^{a+b}(8D)^{|S|}\le(8k)^j$, the coefficient is at least
\[
 \frac1{12\pi d k^3}\frac{j!}{(8k)^j}
 \ge \frac{e^{-4k}}{12\pi k^4}\ge e^{-12k}.
\]
Here $j!\ge(j/e)^j$ and $j\log(k/j)\le k/e$ give the middle inequality. 
The Walsh characters are orthonormal. The Gram matrix of $1,T,T^2,T^3$ on the four $T$ values has smallest eigenvalue at least $2^{-16}$, by its determinant $9/65536$ and trace less than two. This proves the operator bound.

For finite rows, sample independently. Each kernel entry has second moment at most $\E Z_x^2Z_{x'}^2\le3$. Chebyshev and a union bound over $m^2$ entries show that their maximum empirical error is at most $\Lambda_D^{\rm R}/2$ with probability at least $7/8$ at the stated width. Under uniform probability normalization, entrywise error bounds operator norm. The average row squared norm has mean $(d+1)/d\le4/3$ and variance $2(d+1)/(Nd^2)$, so the bound two holds with probability at least $7/8$ after an absolute width threshold. These events have positive joint probability. For such rows use
$c_i=\langle K_N^{-1}r_0,\phi_{a_i}\rangle$; then
$N^{-1}\sum_i c_i^2=\langle r_0,K_N^{-1}r_0\rangle$ gives the claimed representation and coefficient bound.
\end{proof}

\begin{theorem}[Positive ReLU interpolation at near-baseline energy]
For every $\delta>0$ and fixed $A_0>0$ there are $\eta,C,c>0$ and $H_0$ such that, whenever $H\ge H_0$, $1\le D\le\eta H$ and $N\ge e^{CH}$, every function
$r:\mathcal X_D\to[0,A_0]$ is exactly represented by an $H$-hidden-layer width-$N$ ReLU network with
\[
 Q_N(w)\le(1+\delta)H+c.
\]
For every fixed $\kappa>0$ and $t>0$, increasing $C$ gives the following simultaneous finite-scale comparison at every $N\ge e^{C_\kappa H}$ and $0<s\le N^{-1-\kappa}$: a Gaussian mean shift with unchanged native covariance has this mean energy and RMS output error at most $t$ uniformly on the support. Its entropy is exactly $Q_N(w_0)/(2s)$.
\end{theorem}
\begin{proof}
Choose $R>1$ with $R^2\le1+\delta/4$ and set
$r_0=r R^{-(H-2)}$. Use the preceding lemma to implement $r_0$ by a linear combination of first-layer ReLU features. Make every second-layer row equal to $(c_i/N)_i$. Because $r_0\ge0$, the second-layer ReLU leaves it unchanged. In each subsequent hidden layer use the all-constant matrix $R/N$, and use readout weights $1/N$. The resulting output is exactly $r$.

The energy is at most
\[
 2d+2(\Lambda_D^{\rm R})^{-1}A_0^2R^{-2(H-2)}
                  +(H-2)R^2+1.
\]
Choose $\eta$ so small that $2\eta\le\delta/4$ and
$12\eta<2\log R$. The second term is then bounded uniformly in $H$ and $D\le\eta H$, while the first and later-layer terms are at most $(1+\delta/2)H+O(1)$. The finite-row threshold is $e^{O(D)}$, hence at most $e^{CH}$.

For the finite-scale statement write $\tau=\sqrt{Ns}$ and let
$e_\ell=(\E\|h_\ell-h_\ell^0\|^2/N)^{1/2}$ under the native Gaussian perturbation. The first-layer error is at most $\tau\sqrt{2/d}\le\tau$. Its deterministic normalized norm is at most $\sqrt{2d}$. The second mean matrix has uniformly bounded Frobenius norm by the bounded energy term above, so $e_2\le C\sqrt d\,\tau$.
At every later layer the deterministic operator norm is $R$ and the deterministic hidden coordinate is between zero and $A_0$, hence
\[
 e_\ell\le R e_{\ell-1}+\tau(A_0+e_{\ell-1})
       \le(R+1)e_{\ell-1}+A_0\tau.
\]
Thus $e_H\le C\sqrt d\,(R+1)^H\tau$.
The readout perturbation contributes RMS at most
$\sqrt s(A_0+e_H)$, and its mean readout contributes $e_H$.
For $N\ge e^{C_\kappa H}$ and $s\le N^{-1-\kappa}$ this is at most $t$ after enlarging $C_\kappa$. The mean shift preserves covariance, so its Gaussian entropy is exactly its weighted energy divided by $2s$.
\end{proof}

\begin{lemma}[High-temperature control without bounded hidden activations]
Suppose labels are bounded by one, the prior mean prediction is zero, and its empirical mean squared loss is at most $\widehat R(0)+s$. Assume the operator-entropy inequality
\[
 2s\operatorname{KL}(\Pi\|\mu)
 \ge(1-\eta)\E_\Pi Q_{\rm op}-s\mathcal B_\eta
\]
and the output-energy floor
$\|z_w(S)\|_{\rm emp}>t\implies Q_{\rm op}\ge q_H(t)$.
Then every Gaussian-likelihood posterior satisfies
\[
 \|\E_\Pi z_w(S)\|_{\rm emp}
 \le t+C\sqrt{\frac{\lambda^{-1}+s\mathcal B_\eta}{q_H(t)}}
\]
for fixed $\eta<1$ and $s\le1$. The same bound holds with posterior draws clipped to $[-1,1]$ before averaging, and clipping the average cannot increase its norm.
\end{lemma}
\begin{proof}
Jensen gives $F\le\widehat R(0)+s\le2$.
The exact variational identity implies $\E_\Pi\widehat R\le F\le2$, so
$\E_\Pi\|z_w(S)\|_{\rm emp}^2\le6$.
Combining the identity with operator entropy gives
$\E_\Pi Q_{\rm op}\le C(\lambda^{-1}+s\mathcal B_\eta)$.
Split at $\|z_w(S)\|_{\rm emp}\le t$, use Markov on the energy floor, and then Cauchy--Schwarz on the complementary empirical norm. Clipping can only reduce each such norm.
For ReLU on the present equal-norm input support the native prior output variance is
$\frac2d s(\tau^2/2)^H\le s$, so the prior comparison assumption holds exactly.
\end{proof}
